%% kapa: exam=1 round=29 year=2018 period=2 form=A subject=accounting subject_ko=회계학 seq=18 number=58
%% source: Q-Net 감정평가사 기출문제 (cst00309) artlSeq=5202808 file=extracted/1008/5202808/2018년 제29회 감정평가사 1차시험 기출문제/2018년 제29회 감정평가사 1차시험 기출문제/2교시 A형.pdf
%% conversion: kordoc 4.24.51 -> tools/gen1.py (auto); review=proofread-sonnet; flags=none
\begin{question}{58}
충당부채, 우발부채 및 우발자산에 관한 설명으로 옳지 않은 것은?
\begin{choices}
  \item 충당부채는 현재의무이고 이를 이행하기 위하여 경제적 효익이 있는 자원을 유출할 가능성이 높고 해당 금액을 신뢰성 있게 추정할 수 있으므로 부채로 인식한다.
  \item 제품보증이나 이와 비슷한 계약 등 비슷한 의무가 다수 있는 경우에 의무 이행에 필요한 자원의 유출 가능성은 해당 의무 전체를 고려하여 판단한다.
  \item 재무제표는 미래 시점의 예상 재무상태가 아니라 보고기간 말의 재무상태를 표시하는 것이므로, 미래 영업에서 생길 원가는 충당부채로 인식한다.
  \item 손실부담계약은 계약상 의무의 이행에 필요한 회피 불가능 원가가 그 계약에서 받을 것으로 예상되는 경제적 효익을 초과하는 계약을 말한다.
  \item 우발자산은 과거사건으로 생겼으나, 기업이 전적으로 통제할 수는 없는 하나 이상의 불확실한 미래 사건의 발생 여부로만 그 존재 유무를 확인할 수 있는 잠재적 자산을 말한다.
\end{choices}
\end{question}
\answer{3}
